\section{Zusammenfassung}
\label{sec:Zusammenfassung}

In diesem Experiment wurde das Verhalten von dem realen Gas Schwefelhexafluorid untersucht. \\
Durch die Messung von mehreren Isothermen konnten sich die jeweiligen Maxwell-Geraden und eine Dampfdruckkurve bestimmen lassen. 
Daraus lassen sich die kritischen Parameter des Gases ermitteln:

\begin{equation*}
    T_K = (318.6 \pm 2.1) \, K  
\end{equation*}

\begin{equation*}
    p_K = (38,5 \pm 4,0) \, \text{bar}
\end{equation*}

\begin{equation*}
    V_K = (0,32 \pm 0,14) \, \text{cm}^3
\end{equation*}

Mithilfe dieser Parameter konnte die Verdampfungsenergie, die Stoffmenge, die stoffspezifischen Van-der-Waals-Parameter und der kritische Koeffizient bestimmt werden:

\begin{equation*}
    \Lambda = (20.9 \pm 0.50) kJ/mol
\end{equation*}

\begin{equation*}
    \nu = (1.319 \pm 0.014) \cdot 10^{-3} \, mol
\end{equation*}

\begin{equation*}
    a = (0.79 \pm 0.10) \, Pa \frac{m^6}{mol^2}
\end{equation*}

\begin{equation*}
    b = (86.7 \pm 9.6) 10^{-6} \, \frac{m^3}{mol}
\end{equation*}

\begin{equation*}
    K_K = (2.87 \pm 0.12)
\end{equation*}

\noindent Anschließend daran wurden die Messergebnisse mit den theoretischen Vorhersagen verglichen. Dabei zeigte die ideale Gasgleichung nur eine Übereinstimmung bei sehr hohen Volumina und bei kleinen Volumina ist die Abweichung sehr groß.
Die starke Abweichung der idealen Gasgleichung bei kleinen Volumina
ist darauf zurückzuführen, dass sowohl das endliche Eigenvolumen
der Moleküle als auch intermolekulare Wechselwirkungen im Modell
vernachlässigt werden. Insbesondere bei hohen Dichten führen
anziehende Kräfte zu einer Reduktion des Drucks gegenüber dem
idealen Verhalten, während das effektive Ausschlussvolumen
das frei verfügbare Volumen verringert.
Die Van-der-Waals-Gleichung zeigte jedoch eine gute Übereinstimmung mit den gemessenen Werten, versagt jedoch im Zweiphasengebiet, da sie
den Phasenübergang nicht als Koexistenz zweier makroskopischer Phasen
bei konstantem Druck modelliert. Stattdessen liefert sie einen
instabilen Volumenbereich mit negativer Kompressibilität. Die reale
Phasentrennung kann nur durch die Maxwell-Konstruktion angenähert
werden, welche jedoch nicht aus der Zustandsgleichung selbst folgt,
sondern eine zusätzliche thermodynamische Bedingung darstellt.
\\
Die kritische Temperatur wird vergleichsweise gut reproduziert,
da sie primär durch die Form der Isothermen und das Verschwinden
des Koexistenzbereichs bestimmt wird, was experimentell relativ
robust beobachtbar ist. Das kritische Volumen hingegen ergibt sich
aus einer Extrapolation der Maxwell-Endpunkte und ist daher stark
von kleinen Unsicherheiten in Druck und Volumen abhängig.
Bereits geringe Steigungsänderungen führen zu deutlichen
Verschiebungen des Schnittpunkts, sodass \( V_K \) wesentlich
ungenauer bestimmt wird als \( T_K \).\\
Die Bestimmung der Verdampfungsenthalpie ist zusätzlich dadurch
eingeschränkt, dass sie aus der Temperaturabhängigkeit des
Sättigungsdrucks bzw. aus der Clausius-Clapeyron-Beziehung
abgeleitet wird. Diese setzt ein exakt bestimmtes
Phasengleichgewicht voraus. In der Nähe des kritischen Punktes
nimmt jedoch die Latentwärme ab und geht gegen null, wodurch
kleine Messunsicherheiten in Druck und Temperatur zu großen
relativen Fehlern in der berechneten Verdampfungsenergie führen.
\\
Das Molekül SF$_6$ eignet sich besonders gut zur Überprüfung
der Van-der-Waals-Gleichung, da es nahezu kugelsymmetrisch ist
und keine permanenten Dipolmomente besitzt. Dadurch werden
komplexe richtungsabhängige Wechselwirkungen minimiert, sodass
die Annahmen der Van-der-Waals-Theorie — insbesondere die
Beschreibung durch ein effektives Eigenvolumen und eine
mittlere anziehende Wechselwirkung — vergleichsweise gut erfüllt
sind.